\documentclass[11pt]{article}
\usepackage[T1]{fontenc}
\usepackage[margin=1in]{geometry}
\usepackage{amsmath}
\usepackage{amssymb}
\usepackage{enumitem}

\newcommand{\R}{\mathbb{R}}
\newcommand{\N}{\mathbb{N}}
\newcommand{\C}{\mathbb{C}}

\title{Problem Set 4: Sequences, Series, and Integration}
\author{M. Alvarez \\ Math 201, Section B}
\date{October 3, 2025}

\begin{document}
\maketitle

\noindent This problem set covers convergence tests for infinite series, techniques of
integration, and an introduction to linear systems via matrix notation. Show all work;
partial credit is awarded for correct setup even when the final numeric answer is wrong.

\section*{Problem 1}

Let $\{a_n\}_{n=1}^{\infty}$ be the sequence defined by $a_n = \dfrac{2n^2 + 3n - 1}{5n^2 + n}$.

\begin{enumerate}[label=(\alph*)]
  \item Compute $\displaystyle\lim_{n \to \infty} a_n$, showing your reasoning.
  \item Determine whether the series $\displaystyle\sum_{n=1}^{\infty} (a_n - \tfrac{2}{5})$
        converges or diverges, and justify your answer using an appropriate test.
  \item Is the sequence $\{a_n\}$ monotonic for $n \geq 1$? Prove your claim.
\end{enumerate}

\paragraph{Solution.} For part (a), divide numerator and denominator by $n^2$:
\begin{equation}
  \label{eq:limit-an}
  a_n = \frac{2 + 3/n - 1/n^2}{5 + 1/n} \longrightarrow \frac{2}{5} \quad \text{as } n \to \infty.
\end{equation}
By~\eqref{eq:limit-an}, $\lim_{n\to\infty} a_n = \tfrac{2}{5}$, so $a_n - \tfrac{2}{5} \to 0$,
which is only a necessary (not sufficient) condition for convergence of the associated series.
A more careful expansion gives
\[
  a_n - \frac{2}{5} = \frac{5(2n^2+3n-1) - 2(5n^2+n)}{5n(5n+1)} = \frac{13n - 5}{5n(5n+1)},
\]
which behaves like $\dfrac{13}{25n}$ for large $n$, so by the limit comparison test with the
harmonic series $\sum 1/n$, the series in part (b) \emph{diverges}.

\section*{Problem 2}

Evaluate the following integrals.

\begin{enumerate}[label=(\roman*)]
  \item $\displaystyle\int_0^{\pi/2} \sin^3(x)\cos(x)\,dx$
  \item $\displaystyle\int \sqrt{\bigl(4 - x^2\bigr)}\,dx$ using a trigonometric substitution
  \item $\displaystyle\int_1^{\infty} \frac{1}{x^2 + 1}\,dx$, identifying it as an improper integral
\end{enumerate}

\paragraph{Solution.} For (i), let $u = \sin x$, so $du = \cos x\,dx$ and the bounds become
$u=0$ to $u=1$:
\begin{align}
  \int_0^{\pi/2} \sin^3(x)\cos(x)\,dx
    &= \int_0^1 u^3\,du \notag \\
    &= \left[\frac{u^4}{4}\right]_0^1 \notag \\
    &= \frac{1}{4}.
\end{align}

For (ii), substitute $x = 2\sin\theta$, $dx = 2\cos\theta\,d\theta$, so that
$\sqrt{4-x^2} = 2\cos\theta$ and
\[
  \int \sqrt{4-x^2}\,dx = \int 4\cos^2\theta\,d\theta
  = 2\theta + \sin(2\theta) + C
  = 2\arcsin\!\Bigl(\frac{x}{2}\Bigr) + \frac{x\sqrt{4-x^2}}{2} + C.
\]

For (iii), we recognise the antiderivative $\arctan x$ and take a limit:
\begin{align*}
  \int_1^{\infty} \frac{1}{x^2+1}\,dx
  &= \lim_{b \to \infty} \Bigl[\arctan x\Bigr]_1^{b} \\
  &= \lim_{b \to \infty} \arctan b \;-\; \arctan 1 \\
  &= \frac{\pi}{2} - \frac{\pi}{4} \\
  &= \frac{\pi}{4}.
\end{align*}
This integral converges to $\pi/4$, since $\arctan b$ is bounded as $b \to \infty$.

\section*{Problem 3}

Consider the linear system
\[
  \begin{cases}
    2x + y - z = 4 \\
    x - 3y + 2z = -6 \\
    3x + 2y + z = 5
  \end{cases}
\]
which can be written in matrix form $A\mathbf{x} = \mathbf{b}$ with
\[
  A = \begin{pmatrix} 2 & 1 & -1 \\ 1 & -3 & 2 \\ 3 & 2 & 1 \end{pmatrix}, \qquad
  \mathbf{x} = \begin{pmatrix} x \\ y \\ z \end{pmatrix}, \qquad
  \mathbf{b} = \begin{pmatrix} 4 \\ -6 \\ 5 \end{pmatrix}.
\]

\begin{enumerate}[label=(\alph*)]
  \item Compute $\det(A)$ and determine whether $A$ is invertible.
  \item If $A$ is invertible, describe (without full computation) how Cramer's rule would
        be used to solve for $z$.
\end{enumerate}

\paragraph{Solution.} Expanding along the first row,
\[
  \det(A) = 2\bigl[(-3)(1) - (2)(2)\bigr] - 1\bigl[(1)(1) - (2)(3)\bigr]
             + (-1)\bigl[(1)(2) - (-3)(3)\bigr].
\]
Simplifying each bracket,
\[
  \det(A) = 2(-3-4) - 1(1-6) - 1(2+9) = 2(-7) - 1(-5) - 1(11) = -14 + 5 - 11 = -20.
\]
Since $\det(A) = -20 \neq 0$, the matrix $A$ is invertible, and the system has a unique
solution $\mathbf{x} = A^{-1}\mathbf{b}$.

For part (b), Cramer's rule states that
\[
  z = \frac{\det(A_z)}{\det(A)},
\]
where $A_z$ is formed by replacing the third column of $A$ with $\mathbf{b}$.

\section*{Problem 4}

Let $f : \R \to \R$ be given by $f(x) = x^3 - 3x + 1$.

\begin{enumerate}[label=(\alph*)]
  \item Find all critical points of $f$ and classify each as a local maximum, local minimum,
        or neither.
  \item Sketch the general shape of the graph, noting concavity changes.
\end{enumerate}

\paragraph{Solution.} We have $f'(x) = 3x^2 - 3 = 3(x-1)(x+1)$, so the critical points are
$x = 1$ and $x = -1$. Since $f''(x) = 6x$, we get $f''(1) = 6 > 0$ (local minimum) and
$f''(-1) = -6 < 0$ (local maximum). The inflection point occurs where $f''(x) = 0$, i.e.\
at $x = 0$, where concavity switches from concave-down (for $x<0$) to concave-up
(for $x>0$). A useful sanity check on the algebra above is the bound
\[
  \left|\, f(x) - f(1) \,\right| \;\le\; \bigl|\, x^3 - 3x \,\bigr| + 2 \qquad \text{for } x \in [-2,2],
\]
which follows directly from the triangle inequality applied term by term.

It is also worth recording the values of $f$ and $f'$ at the critical points side by side,
since this pair of columns is exactly what we will plot in the next section:
\begin{align*}
  f(-1) &= 3, & f'(-1) &= 0, \\
  f(1)  &= -1, & f'(1)  &= 0.
\end{align*}
Finally, the closed-form expression for the local extrema in terms of a generic cubic
$g(x) = x^3 + px + q$ can be written compactly as
\[
  x_{\pm} = \left( \pm\sqrt{\frac{-p}{3}} \right),
\]
which for $p = -3$ recovers $x_{\pm} = \pm 1$ as found above.

\end{document}
